\documentclass[10pt,a4paper]{amsart}
\usepackage[margin=19mm]{geometry}
\usepackage{amsmath,amssymb,mathtools}
\usepackage[scr=rsfs,cal=euler]{mathalfa}
\usepackage{xfrac}
\usepackage[unicode,hidelinks,bookmarks=false]{hyperref}
\newcommand{\ie}{i.e.\ }
\newcommand{\cf}{cf.\ }
\newcommand{\resp}{respectively\ }
\newcommand{\Subsection}{\subsection}
\providecommand{\nobreakdash}{\nobreak\hbox{-}}
\setlength{\emergencystretch}{2em}
\multlinegap0pt
\usepackage{float}
\usepackage{tcolorbox}
\usepackage{bm}
\usepackage{tikz}
\usepackage{xcolor}
\newtheorem{proposition}{Proposition}
\title{A domain decomposition strategy for natural imposition of mixed boundary conditions in port-Hamiltonian systems}
\author{S.D.M. de Jong, A. Brugnoli, R. Rashad, Y. Zhang, S. Stramigioli}
\date{Source: arXiv:2501.06107v4 (CC BY 4.0)}
\begin{document}
\maketitle

\noindent\emph{Source and licence.} A domain decomposition strategy for natural imposition of mixed boundary conditions in port-Hamiltonian systems. Version arXiv:2501.06107v4 (CC BY 4.0), distributed by arXiv under the Creative Commons Attribution 4.0 International licence, http://creativecommons.org/licenses/by/4.0/, as recorded in the arXiv licence metadata for this exact version and shown on the abstract page https://arxiv.org/abs/2501.06107v4. Full licence notice: \texttt{licenses/arxiv-2501.06107-CC-BY-4.0.txt}.\par\smallskip

\noindent\textit{Editorial scope.} Counted unit: the article's proposition proposition (source0.tex lines 1029--1031). The article's complete original proof of it is delivered with it (source0.tex lines 1032--1062, one \texttt{proof} environment of 31 lines). No mathematics is written, repaired or re-derived: the statement and the proof are the article's own lines, in the article's own notation, and no retained line is substituted or re-flowed.  No further source-local context is needed: every cross-reference of every retained line resolves inside the two delivered spans.  Nothing was omitted from any retained span, so the package carries no omission record.  No retained line contains a citation, so the wrapper carries no bibliography and no bibliography entry.  No retained line contains a figure environment, an image-inclusion command or drawing code, so no image is delivered and the package declares no asset dependency.  Editorial formatting only, in addition: an AMS \texttt{amsart} preamble that restates the article's own declarations and macros used by the retained lines, namely the environments \texttt{proposition} on separate counters, together with the standard packages those lines need.  The retained environments print in this document as Proposition 1 in order of appearance, whereas the article prints the same items with the numbering of its own full document.  The complete native source of the article as distributed in the arXiv e-print for this exact version is bundled unchanged as \texttt{sources/171448/source0.tex}.
\begin{proposition}
The implicit midpoint scheme applied to a linear Poisson system is a Poisson map.
\end{proposition}

\begin{proof}
    Using a time rescaling, we set $\Delta t/2 = 1$. 
    The application of the midpoint rule leads to the recursion 
    \begin{equation}\label{eq:discrete_flow}
        {\mathbf{e}}^{n+1} = \mathrm{Cay}({\mathbf{J}}) {\mathbf{e}}^{n}, \qquad  \mathrm{Cay}({\mathbf{J}}) := (\mathbf{I} -{\mathbf{J}})^{-1}(\mathbf{I} + {\mathbf{J}}).
    \end{equation}
    For the discrete flow \eqref{eq:discrete_flow} to be a Poisson map, it must hold
    $$
    \mathrm{Cay}({\mathbf{J}}) \; \mathbf{J} \; \mathrm{Cay}({\mathbf{J}})^\top = \mathbf{J}.
    $$
    By exploiting the property $\mathbf{J} = -\mathbf{J}^\top$, the term $\mathrm{Cay}({\mathbf{J}})^\top$ gives
    $$
    \mathrm{Cay}({\mathbf{J}})^\top = (\mathbf{I} -{\mathbf{J}})(\mathbf{I} + {\mathbf{J}})^{-1}.
    $$
    So the discrete flow can be rewritten as
    $$
    \mathrm{Cay}({\mathbf{J}}) \; \mathbf{J} \; \mathrm{Cay}({\mathbf{J}})^\top =  (\mathbf{I} -{\mathbf{J}})^{-1}(\mathbf{I} + {\mathbf{J}}) \mathbf{J} (\mathbf{I} -{\mathbf{J}})(\mathbf{I} + {\mathbf{J}})^{-1}.
    $$
    The following commuting properties holds
    \begin{equation}\label{eq:commutation_matrices}
        \begin{aligned}
            \mathbf{J}(\mathbf{I} + \mathbf{J}) = (\mathbf{I} + \mathbf{J}) \mathbf{J}, \qquad
            \mathbf{J}(\mathbf{I} - \mathbf{J}) = (\mathbf{I} - \mathbf{J}) \mathbf{J}, \qquad
            (\mathbf{I} + \mathbf{J})(\mathbf{I} - \mathbf{J}) = (\mathbf{I} - \mathbf{J})(\mathbf{I} + \mathbf{J}).            
        \end{aligned}
    \end{equation}
    Using these relations, it is obtained
    $$
    (\mathbf{I} -{\mathbf{J}})^{-1}(\mathbf{I} + {\mathbf{J}}) \mathbf{J} (\mathbf{I} -{\mathbf{J}})(\mathbf{I} + {\mathbf{J}})^{-1} = \mathbf{J}.
    $$
\end{proof}

\end{document}
